%actually  My material: your own input. This is what an AI may draft from, and it is
% graded, because it is your thinking. Keep earlier entries; add new ones.
%   F1  bullet notes for About my research (if AI drafted it)
%   F3  Heilmeier bullet answers; one sentence in Booth's template:
%       "I am studying X because I want to find out Y in order to help
%        readers understand Z."
%   F5  citeable notes (2 to 3 per paper, each with its citation)
%   F7  synthesis matrix (every cell marked verified) and argument outline
%       (3 to 5 themes, the claim of each, supporting papers, the gap)
%   F9  your own first version of the focus statement

\section{My Material}

\subsection{Notes for About My Research (F1)}

\textit{No AI drafted this section of work}

\subsection{Heilmeier Bullets and Research Question (F3)}

\textbf{Question 1:}
\begin{itemize}
    \item Build a test system and control hardware for a satellite propulsion system.
    \item Ensure that the system fires well and handles all potential errors and firing cases
    \item Develop flight-ready hardware and software which can handle the control aspects of the propulsion system.
    \item Find the necessary metrics for good data on the system itself.
\end{itemize}

\textbf{Question 2:}
\begin{itemize}
    \item Power estimates for cubesats: <10W @ 3U
    \item Water vaporiztion: Requires significant energy conumption for liquid-gas phase change
\end{itemize}

\textbf{Question 3:}
\begin{itemize}
    \item Steamjet has already done something similar
    \item Thermal battery: use less instantaneous power to put less of a drag on the system itself.
    \item Nothing new in terms of electronics: combining existing technologies which can be purchased commercially
    \item Already have some control hardware; the rest of the work is data collection, validation, and software.
\end{itemize}

\textbf{Question 4:}
\begin{itemize}
    \item Much more efficient way to provide propulsion to a cubesat
    \item Suited for low volume, low power demand systems
\end{itemize}

\textbf{Question 5:}
\begin{itemize}
    \item PCB delays and hardware issues (mistakes are costly in terms of time) -- Not much I can't fix by hand!
    \item Software has many hurdles
    \item I've created a similar flight algorithm for another class
    \item Vacuum chamber hardware is located out of state.
\end{itemize}

\textbf{Question 6:}
\begin{itemize}
    \item Probably less than a few hundred dollars
    \item Say \$1000 for good measure
    \item Funding is coming from AMSAT
\end{itemize}

\textbf{Question 7:}
\begin{itemize}
    \item around 1 year for the entire project
    \item already 5 months in.
\end{itemize}

\textbf{Question 8:}
\begin{itemize}
    \item Month 3: Control system hardware should be completed
    \item Month 6: ready to fire thruster (plus some form of fast data transfer)
    \item Orbit determination from scratch
    \item Flight-ready hardware design before presentation is commplete
\end{itemize}

I am studying cubesat resistojet propulsion systems because I want to find out if power consumption, size, and mass can be improved in order to help readers utilize the learnings from this project to integrate similar systems on their spacecraft.

\subsection{Citeable Notes (F6)}

\subsection{Finite-time attitude set-point tracking for thrust-vectoring spacecraft rendezvous}

\textbf{Citable Notes} 

\begin{enumerate}
    \item \cite{zhangFinitetimeAttitudeSetpoint2020} shows that the optimal attitude and $\Delta v$ parameters can be independently solved for with different methods. Given the much more complex orbital solution, the attitude pointing solution becomes a closed-form optimization via the multiplier method. 
    \item Providing an improved particle swarm optimization (PSO) algorithm, \cite{zhangFinitetimeAttitudeSetpoint2020} advances PSO techniques for both $\Delta v$ and firing time to meet some targeted rendezvous point. One contstraint on the utility of this method is that PSO, no matter how thorough, does not guarantee a global minimum on the search space.
\end{enumerate} 

\textbf{My Summary}

This work claims that a spacecraft rendezvous system can be created which uses a heuristic optimization method (Particle swarm optimization) which allows a separate pointing solution to be solved in closed form. The simulated spacecraft is a single-thruster satellite which is attempting to meet with a target satellite in a non-concentric orbit. Given the initial position of the first satellite, what's the best time and $\Delta v$ requirements for a series of burns to get to meet with the satellite. This is solved for with an improved particle swarm optimization, a type of heuristic optimization where a local minima is found. The authors improve on this algorithm with a custom implementation that reduces $\Delta v$ requirements. In addition, this paper shows that once the burn timing and $\Delta v$ are solved for, the pointing solution can be solved in closed form using the method of multipliers. They solve this with a lyapunov function.

This optimization method is demonstrated in a four-burn orbital transfer which successfully meets with a target spacecraft. Their PSO method is only marginally better than the existing PSO, and might not be worth using in many cases. Their work also gives a framework, but only tests it over a single orbital maneuver; it does not show the effects of not meeting their assumptions or over different orbital conditions. Cubesat propulsion might have a difficult time with this method as their thrusters are much less powerful, which adds additional constraints that may make PSO fail. They don't compare against a baseline method other than their improved PSO and the standard PSO algorithms, as other methods for orbital transfers might struggle to meet these goals. Their measurement is based only on how much $\Delta v$ was used (i.e., fuel consumed for the maneuver).

As far as supporting work, code and data are not available, but the paper does include their solution parameters to remake their existing code. If I were to use this in my own work, the pointing solution would need to be updated to account for the maximum-power path as well as the pointing solution. For a near-circular orbit, the timing and $\Delta v$ would already be determined beforehand asmy work doesn't include a dimension to track a target spacecraft. Adding the dimension of the spacecraft power generation may move that solution out of its closed form, since speed isn't the only thing we need.

\subsection{Self-Adaptive Energy-Aware Power Management for CubeSat Electrical Power Systems Under Dynamic Orbital Conditions}

\textbf{Citable Notes} 

\begin{enumerate}
    \item \cite{akarneSelfAdaptiveEnergyAwarePower2026} demonstrates a simulated sparrow-search algorithm (SSA) to tune MPC parameters to adaptively correct the time horizon and cost function associated with the MPC algorithm.
    \item Generating a model of a spacecraft power system, \cite{akarneSelfAdaptiveEnergyAwarePower2026} improves upon battery state of health (SOH) over a two-year period using an adaptive MPC algorithm to manage power usage over varying orbital conditions. The model could be extended to handle a thermal battery in parallel with a electric battery.
    \item Using an adaptively-tuned MPC algorithm, \cite{akarneSelfAdaptiveEnergyAwarePower2026} reports an increase in on-demand bus power. Its algorithm is robust to variations in orbital conditions and manages spacecraft power by predicting future loading and power generation.
\end{enumerate} 

\textbf{My Summary} 

This work provides a framework for improving model-predictive control algorithms for spacecraft power management. Model-predictive control utilizes a time horizon and cost weight parameters to determine what features get most optimized for over time. This work's innovation is to handle adaptively tuning this MCP optimization strategy with a sparrow-search algorithm (SSA) that automatically tunes the MPC parameters to take into account the power generation, load needs, and future orbital states. The SSA-MPC algorithm is compared with a fixed-parameter MPC as a baseline over a single orbit.

The SSA-MPC also takes into account external spacecraft parameters like termperatures and eclipse times. Using a model of battery state of health (SOH) degredation over time, the battery is shown to have dropped 4\% less in SOH over a fixed-parameter MCP. It also shows that the margin of residual power leftover during the eclipse period can be much greater in different orbital conditions. Overall, their claim that their method is more energy-efficient than existing methods is certainly shown.

They don't take into account orbital variation well and don't consider loads with probabilistic models of load consumption. For example, a communication payload does not have the ability to be scheduled like other loads. Packets arrive randomly and must be processed and responded to, which is a significant load on the satellite. The power this operation consumes is therefore based on the probability of packets arriving in the first place. In addition, they do a monte-carlo analysis over variations in an orbit, but don't consider a change in orbit like a thruster would take. What happens if orbital parameters change completely? This seems like a good way to handle that, but they don't show it. 

This relates exactly to my work as the satellite with the propulsion system on it adds a thermal battery for energy storage which has its own heath degredation over time. To prepare for a firing, the load would have to adapt to the additional thruster payload, which is how we might be able to keep the satellite running even with the extra payload of the thruster.

\subsection{Spacecraft Attitude Estimation from Vector Measurements Using Particle Filter} 
\textbf{Citable Notes} 

\begin{itemize}
    \item \cite{xue-yuanjiangSpacecraftAttitudeEstimation2005} constructs a particle filter (PF) implementation of an attitude estimator using gyro and three-axis magnetometer (TAM) measurements, reporting an increase in convergence time between a UKF and PF implementation of the same algorithm.
    \item Focusing on convergence between the UKF and PF implementation, \cite{xue-yuanjiangSpacecraftAttitudeEstimation2005} demonstrates that the PF outperforms the UKF in terms of convergence time for attitude estimation. Unfortunately, this work gives no indication of why the PF outperforms the UKF in terms of convergence time or the absolute attitude error. 
\end{itemize}

\textbf{My Summary} 

This work covers a particle filter estimator for a satellite with a gyro and three-axis magnetometer. The particle filter is compared against a UKF implementation using the same model. The authors compare the particle filter representation and the UKF in terms of convergence time. This work does not compare the absolute accuracy of either method even though all data is based on simulation. It is shown that the particle filter convergence time is much shorter than the UKF when the initial uncertainties are very high. Additionally, the authors don't do anything to characterize approximate cost of either methods's calculations.

This work relates to mine as a backup orbit determination is planned for the rest of this work. One thing that this work adds is a model of a tumbling spacecraft. I've done similar work on the particle filter in the past, but tumbling spacecraft which are also controlled (i.e., detumbling) can be difficult to model. This work uses generalized Rodrigues parameters (GRPs) with a quaternion to model the rotation. The GRP is a parameterization of the quaternion which decreases the dimension of the measurement vector and is well-suited for a particle filter representation.

The authors provide no code and not enough information to reproduce their work. There are some sources which are cited as doing similar things as this work which may act as a better reference than this one. As this is a conference paper, they may have been limited by page count, so didn't get all their testing or findings out with the paper. 

\subsection{AI Report Summary} 

I had AI create reports of all three papers since it wasn't that much extra work. I was really surprised at how much it got right compared with what the lecture had said for this week. My strategy was to read and take notes, then have the AI convert the paper to a text file with the PDF conversion tools, read it, answer the questions I put above, and summarize what it said. This way, the paper would be copied directly from the source instead of the model trying to pull it from its weights. This worked well: there were very few mistakes and it was pretty much always able to answer my questions. I think the biggest thing it misses is the connection to my work. For example, in the first paper I read, it did a lot of qualification and explanation of the particle swarm optimization. While that was interesting, I don't think it's quite what I needed because my work isn't going to use that technique at all. I was only looking for the pointing part and it doesn't say anything about how to meet the goals outlined in my prospectus. Also, coming up with ideas even given context of what my project is usually is not worthwhile from AI, in my opinion. 

There were definitely things it got wrong or misunderstood. In the second paper I read, it added an entire section of notes on the antenna gains, doppler fading, and pass geometry. This actually has very little to do with the paper. I suppose it's not technically wrong, but completely irrelevant and would have confused me had I not read the paper. For the last paper I read, I had a question about why they qualify something when both orbits converged in a very short amount of time. This was based on one of the plots that they show in the paper. The AI claimed that the convergence was 3 hours for one and 1 hour for the other. I have no idea where this came from, as the text never discusses this numerically. The AI was not able to read the plot and just made something up.

I think that the summaries were not helpful in comparison to just reading the paper. They essentially wrap everything I didn't understand into a more dense and less readable form. They have even less context than the paper does and I don't entirely trust it since I don't understand it well. If there was a mistake, I wouldn't necessarily catch it in something I'm unfamiliar with. I think the best use of it is to answer specific question from my notes on the paper. Asking AI to help once the paper is somewhat understood or having it recreate the code for some result might be helpful. At the end of the day, I would want to thoroughly understand what the paper says rather than repeat what an AI says.


\subsection{Synthesis Matrix and Argument Outline (F7)}

\textbf{Topic: Power generation can be scheduled to handle changing orbital conditions (no AI used)}

\begin{longtable}{@{}p{0.08\linewidth}p{0.24\linewidth}p{0.13\linewidth}p{0.12\linewidth}p{0.15\linewidth}p{0.15\linewidth}@{}}
  \toprule
    Paper & Approach & Baseline & Metric & Result & Limitation \\
  \midrule
  \endhead
    \cite{husseinOptimizedLoadSchedulingAlgorithm2023} & 
    \begin{minipage}[t]{\linewidth}
    \begin{itemize}[nosep, left-margin=1em]
        \item Fixed scheduling: power is optimized over a fixed time horizon.
        \item BER and data link: maximize the power to the data link when channel supports it
        \item Place different loads into different categories by time, magnitude, or time and magnitude\strut
    \end{itemize}
    \end{minipage} &
    Compared to timing based communication approach &
    BER, energy consumed, data rate&
    significantly reduced BER over naive approach &
    Model of other satellite loads and power generation is extremely simple \\
    \addlinespace

    \cite{akarneSelfAdaptiveEnergyAwarePower2026} & 
    \begin{minipage}[t]{\linewidth}
    \begin{itemize}[nosep, left-margin=1em]
        \item Modifed MPC approach where parameters are adaptively tuned given what is known about the orbit.
        \item Battery SOH, power generation, and external loads are significantly more realistic compared with other studies.
        \item Manages battery storage SOC and temperature
        \item Accounts for future power generation and loading\strut
    \end{itemize}
    \end{minipage} &
    Compared to bare MPC approach &
    Battery SOC, Battery SOH, residual power & 
    4\% improvement on SOH over 2-year period and increased residual power available in eclipse & 
    Doesn't actually simulate over multiple orbits; no estimators of system parameters from noisy measurements \\
    \addlinespace
    
    \cite{doff-sottaSpacecraftEnergyManagement2022} & 
    \begin{minipage}[t]{\linewidth}
    \begin{itemize}[nosep, left-margin=1em]
        \item Introduces a fully-convex optimization method for a GEO satellite power system 
        \item Fully open-loop: doesn't really control for anything except battery SOC and temperature
        \item Satellite model is much larger than standard cubesat formfactor (on order of 35 kW power allocation)
        \item Pushes towards convexity by assuming that heat generation and external factors can all be assumed beforehand; no measurements of real parameters
        \item Optimization is meant to maximize power allocated to secondary tasks while taking into account necessary loads.\strut
    \end{itemize}
    \end{minipage} &
    Compared to bare MPC approach &
    Battery stored energy and computation time&
    Improvement on battery storage over singular orbit &
    Doesn't take into account any real measurements in the satellite; I don't really trust that the measurements the paper presents would be accurate in space. \\
  \bottomrule
\end{longtable}

\textbf{Topic: Pointing can be optimized for time or power (no AI used)} 
\begin{longtable}{@{}p{0.08\linewidth}p{0.24\linewidth}p{0.13\linewidth}p{0.12\linewidth}p{0.15\linewidth}p{0.15\linewidth}@{}}
  \toprule
    Paper & Approach & Baseline & Metric & Result & Limitation \\
  \midrule
  \endhead
    \cite{zhangFinitetimeAttitudeSetpoint2020} & 
    \begin{minipage}[t]{\linewidth}
    \begin{itemize}[nosep, left-margin=1em]
        \item Optimize first for $\Delta v$ then the Pointing optimization can be done in closed form.
        \item Utilize PSO instead of a better optimization technique to handle the multi-transfer, high-dimensional problem.
        \item Assume a large satellite that has the capability of making a multi-burn orbital transfer.\strut
    \end{itemize}
    \end{minipage} &
    Not compared against other techniques &
    Time and accuracy of the maneuver (did the chaser meet the rendezvous?) & 
    Successful implemention over a single multi-transfer maneuver&  
    No actual optimality; Only solve for optimized-time pointing maneuver \\
    \addlinespace
    \cite{kristiansenEnergyOptimalAttitude2024} & 
    \begin{minipage}[t]{\linewidth}
    \begin{itemize}[nosep, left-margin=1em]
        \item Created a power-based cost function to minimize the control objective while maximizing energy generation (pointing)
        \item Simulated control technique over a multi-objective orbit and calculated power generation.
        \item Included eclipse period in optimization formulation solved with IPOPT and CasADi.\strut
   \end{itemize}
   \end{minipage} &
    Proportional-Derivative controller doing the same task &
    Final pointing error and increase in energy generation & 
    39\% increase in energy generation while meeting pointing spec. &  
    Not feasible for real-time applications; not well-suited for propulsive maneuvers \\
  \bottomrule
\end{longtable}

\textbf{Topic: Particle filters can be used to estimate attitude of a spacecraft (no AI used)} 
\begin{longtable}{@{}p{0.08\linewidth}p{0.24\linewidth}p{0.13\linewidth}p{0.12\linewidth}p{0.15\linewidth}p{0.15\linewidth}@{}}
  \toprule
    Paper & Approach & Baseline & Metric & Result & Limitation \\
  \midrule
  \endhead
    \cite{xue-yuanjiangSpacecraftAttitudeEstimation2005} & 
    \begin{minipage}[t]{\linewidth}
    \begin{itemize}[nosep, left-margin=1em]
        \item Utilize Gyro and TAM as estimator inputs for state estimation of the attitude of the satellite.
        \item Develop model of spacecraft attitude which handles tumbling and any rotation.\strut
    \end{itemize}
    \end{minipage} &
    Unscented Kalman Filter using the same model & 
    Convergence time and initial attitude covariance &
    Inproved convergence time for high-uncertainty systems &
    No baseline error calculation against absolute attitude error; doesn't account for propulsion \\
  \bottomrule
\end{longtable}

\subsection{Argument Outline} 
\textbf{Themes:}
\begin{enumerate}
    \item Power scheduling is a separate optimization task which can be scheduled to adapt to changing orbital conditions.
    \begin{itemize}
        \item Power scheduling tasks are structured by the problem attempting to be solved.
        \item Maximize power to load, maximize available power at any time, minimize battery degredation over time, and minimize unused power.
        \item \cite{husseinOptimizedLoadSchedulingAlgorithm2023} demonstrates a method of load scheduling to maximize power to a communication link when the link is most likely to be successful (BER is minimized). 
        \item Although successfully minimizing BER, the method employed here is simplistic given changing mission and orbital conditions. 
        \item Their model never takes into account pointing requirements, designating everything not involved with the communication link as a fixed load. 
        \item This also assumes that the communication payload is fixed and must be sent at intervals which maximize LOS (not necessarily true in this case, as operations happen at a random and constant rate with this mission)
        \item \cite{doff-sottaSpacecraftEnergyManagement2022} implements a fully-confex formulation of the power problem for a much larger spacecraft, and shows how certain constraints could be made into convex constraints, thereby making the problem much easier and faster to solve.
        \item This work assumes that generated power will generate more energy than can be stored in batteries onboard.  
        \item In addition, no pointing constraints, feedback from the system (entirely open-loop control), or control over indiviual loads is considered.
        \item This work requires that the slewing control be taken into account, as well as maximizing the power to a specific load; still understanding that some constraints could be made convex is useful
    \end{itemize}
    \item This mission specifically is well-suited for power scheduling since the goal is to optimize over many load conditions.
    \begin{itemize}
        \item This mission's goal is to maximize available power up to the burn from the propulsion system.
        \item The propulsion system should be fed all available power to phase the PCM in as little time as possible.
        \item From experiments, it will take at least 2 hours at maximum power to get to temperture.
        \item \cite{akarneSelfAdaptiveEnergyAwarePower2026} demonstrates a method of managing multiple loads at once while maximimizing battery state of charge and minimizing battery wear over time.
        \item This work creates model predictive control algorithm which iteratively searches for model optimization parameters using a sparrow search algorithm
        \item It is able to account for orbital changes and handle multiple, time and magnitude-varying loads.
        \item There is no formulation of a probabilistic model of radio transmission for AMSAT needs (for example, packets aren't dependent on specific ground stations, but are coming as discrete bursts as people send packets all over the world).
        \item The gap comes as they don't simulate over an orbital change or plan for the specific situation that our satellite will be in. This formulation could be expanded to include the load of the thermal battery and PCM degredation included with the thermal battery.
    \end{itemize}
    \item The optimal pointing solution has a closed-form solution when time is the objective, but requires a non-linear optimization schedule when maximum power is desired.
    \begin{itemize}
        \item The thruster is aligned with a single axis of the spacecraft, and the pointing direction necessary for thrust will require slewing into a certain position.
        \item In addition, there may be communication tasks which require some amount of slewing to get into position.
        \item \cite{zhangFinitetimeAttitudeSetpoint2020} demonstrates that once a set of burns for a maneuver is known (solved for there using a multi-objective PSO) the power-minimizing pointing objective can be found using a closed-form solution
        \item Their work is for a much larger satellite (gap) and does not require that a satellite maintains optimal pointing towards the sun at all times.
        \item Also assumes that we require a complex PSO solution to solve for the burn $\Delta v$ and times, which may be true if we need to make multiple burns to meet the orbital transfer goals that would be required.
        \item \cite{kristiansenEnergyOptimalAttitude2024} builds off existing work to form an non-linear optimization for maximizing power when slewing to a set of positions for both communication and payload tasks
        \item Their formulation is much more complex and utilizes a non-linear solver(IPOPT) to handle the optimal slewing, although they generate $\approx$ 40\% more energy over the entire maneuver while meeting their much more complex pointing requirements than this mission requires.
        \item They generate a good model of reaction wheel and magnetorquer dynamics which could be used in other formulations where we're trying to also balance loading with the propulsion system.
        \item This model lacks any type of pointing for propulsion: a highly-weighted objective should happen at the highest $v$ point to slew the satellite into the direction of motion of the bird.
    \end{itemize}
    \item A sequential PF to UKF algorithm is well-suited to potentially detumble after firing to reduce computation time.
    \begin{itemize}
        \item \cite{xue-yuanjiangSpacecraftAttitudeEstimation2005} utilizes a particle filter to estimate the attitude of the satellite given gyroscoping and three-axis magnetometer (TAM) measurements. 
        \item Their work is more of a demonstration than analytical proofs of certain features of the particle filter; they show that the particle filter converges faster for large initial estimation errors when compared with a UKF implementation of the same algorithm.
        \item The gap in their work is that they never check against absolute error, only relative to the error of the UKF, and have no explanation of some of the features of the particle filter.
        \item This work proposes to handle a more rigorously tested algorithm that tests their conclusions over more cases, and adds in orbital estimation or some kind of hybrid estimation utilizing GPS measurements as well.
    \end{itemize}
\end{enumerate}


\begin{framed}
    \textbf{AI Text Generation Prompt: }
    I've written up an outline for a literature review section in @appendices/
    my\_material.tex (see the bottom of the page). I would like you to write a 1.5-2 page literature review based on that outline, using exactly the citations that I use in the materials section in @sections/literature\_review.tex. Please create a copy of your writing and place it in sections/litreview\_ai\_v0.tex as well as adding \texttt{\\trackededits{litreview\_ai\_v0}} in @appendices/tracked\_edits.tex and an addition in @appendices/ai\_use\_log.tex for what you did. You can find my notes on each one of these papers at ~/.local/share/lx/notes/20261002-propulsion-notes-runnign.tex and you can find the papers in ~/Library/Papers/Thesis/ in directories associated with their topics. Feel free to convert these to text to understand what's being said in each of them. Try to stick to claims in the actual material instead of running off things you remember; if you can't find a source for it in my notes or one of the papers, you probably should leave it as a question. Feel free to leave comments as questions to me or things that you think I might have gotten wrong.
\end{framed}

\begin{figure}[H]
    \centering
    \includegraphics[width=0.75\textwidth]{ResearchConnections.png}
    \caption{ResearchRabbit search on core papers colored by red (pointing), purple (power), blue (orbit/attitude), and white (not in collection). I have added some of these papers to my collection since creating this graph.}
\end{figure}

\subsection{First Version of My Focus Statement (F9)}
\textit{[Your own version, before any AI help.]}
